\chapter{Stratégies de manipulation et manière dont elles sont détectées}\label{ch:s2:manipulative-strategies-and-how-they-are-caught}

Un trader plaçait de gros ordres qu'il ne voulait jamais exécuter, les annulait en une fraction de seconde et négociait du côté opposé du mouvement de cours
qu'ils provoquaient. Dans l'affaire United States v. Coscia, la Seventh Circuit a confirmé sa condamnation, la première fondée sur la disposition \omterm{def:s2:manipulative-strategies-and-how-they-are-caught:spoofing}{anti-spoofing} du Commodity Exchange Act.
Les preuves figuraient dans le journal des ordres~: 24\,814 gros ordres sur trois mois de 2011, dont 0.5~\% ont été exécutés, face à de petits ordres du côté opposé dont
environ 52~\% ont été exécutés. Ce chapitre décrit chaque pratique de manipulation uniquement dans la mesure où le dossier public le permet~: son fonctionnement, les raisons pour lesquelles elle est illégale ou
contestée, l'affaire qui l'a établie et le détecteur qui la repère. Sur des jours de compte synthétiques, un détecteur qui combine les deux caractéristiques du
dossier Coscia détecte 77~\% du \omterm{def:s2:manipulative-strategies-and-how-they-are-caught:spoofing}{spoofing} injecté à un taux de faux positifs de 1~\%. Le ratio ordres-transactions, souvent cité comme signal de \omterm{def:s2:manipulative-strategies-and-how-they-are-caught:spoofing}{spoofing}, n'en détecte aucun~:
les teneurs de marché annulent davantage que les spoofeurs. La construction repose sur \texttt{firm.surveil}.

\section{Spoofing et layering}

\begin{definition}[Spoofing]\label{def:s2:manipulative-strategies-and-how-they-are-caught:spoofing}
Le \emph{spoofing}\index{spoofing} consiste à passer un ordre d'achat ou de vente dans l'intention de l'annuler avant son exécution, de façon à créer une
fausse impression d'offre ou de demande et à négocier contre le mouvement de cours qu'il provoque.
\end{definition}

\begin{definition}[Layering]\label{def:s2:manipulative-strategies-and-how-they-are-caught:layering}
Le \emph{layering}\index{layering} est un \omterm{def:s2:manipulative-strategies-and-how-they-are-caught:spoofing}{spoofing} qui consiste à placer plusieurs ordres à des niveaux de prix successifs d'un même côté du carnet, tenus à l'écart du meilleur
prix, de sorte qu'ils soient vus mais peu susceptibles d'être exécutés.
\end{definition}

\textbf{Fonctionnement.} Selon l'exposé de l'opinion Coscia, un programme plaçait un petit ordre d'un côté et de gros ordres de l'autre. Les gros ordres
suggéraient une offre ou une demande qui n'existait pas, et poussaient le cours vers le petit ordre, qui était exécuté. Les gros ordres étaient alors annulés,
et le processus se déroulait en sens inverse. Chaque cycle durait environ deux tiers de seconde. Il a été répété pour plus de 450\,000 gros ordres et a rapporté
1.4~millions de dollars. La CFTC a décrit la version de Sarao sur le contrat à terme E-mini S\&P 500~: à partir de 2009, une plateforme de trading modifiée plaçait en \omterm{def:s2:manipulative-strategies-and-how-they-are-caught:layering}{layering} quatre à six ordres de vente exceptionnellement
importants, séparés d'un niveau de prix et situés à au moins trois ou quatre niveaux du meilleur cours vendeur, et la plupart étaient annulés sans avoir été exécutés.

\textbf{Pourquoi c'est illégal.} Le Dodd-Frank Act de 2010 a ajouté au Commodity Exchange Act une interdiction du \omterm{def:s2:manipulative-strategies-and-how-they-are-caught:spoofing}{spoofing}, défini comme le fait de placer un ordre d'achat ou de vente dans l'intention
de l'annuler avant exécution. Le Market Abuse Regulation de l'Union européenne range parmi les manipulations les ordres qui donnent des signaux faux ou trompeurs, y compris
les ordres placés puis annulés par des algorithmes, et cite comme indicateur les ordres qui modifient le meilleur cours acheteur ou vendeur ou la représentation du carnet et sont
retirés avant d'être exécutés. L'infraction tient à l'intention, non à l'annulation~: la plupart des ordres sont annulés, et Coscia soutenait que les traders à haute fréquence
annulent 98~\% des leurs.

\textbf{Les affaires de répression.} Coscia a été reconnu coupable sur tous les chefs d'accusation, condamné à 36 mois d'emprisonnement et a perdu son appel en août 2017. La cour a retenu l'intention
à partir de la conception du programme~: celui-ci annulait les gros ordres après un délai fixé, si les petits ordres étaient exécutés ou si un gros ordre unique était exécuté. Sarao a été poursuivi
par la CFTC et le ministère de la Justice en avril 2015. La CFTC a allégué qu'il avait utilisé le programme de \omterm{def:s2:manipulative-strategies-and-how-they-are-caught:layering}{layering} pendant plus de 400 jours de négociation et avait gagné plus de
40~millions de dollars, et avait exercé pendant deux heures, avant le flash crash du 6 mai 2010, une pression vendeuse persistante de près de 200~millions de dollars. En novembre 2016, il
a accepté une sanction de 25.7~millions de dollars et la restitution de 12.9~millions de dollars.

\begin{dated}{2026-09}{Le droit}\label{dat:s2:manipulative-strategies-and-how-they-are-caught:law}
Aux États-Unis, 7~U.S.C.~6c(a)(5)(C) rend illicite toute négociation qui relève, présente le caractère ou est communément connue dans le métier comme du \omterm{def:s2:manipulative-strategies-and-how-they-are-caught:spoofing}{spoofing}, défini
comme le fait de placer un ordre d'achat ou de vente dans l'intention de l'annuler avant exécution. Dans l'Union européenne, l'article~12 du Market Abuse Regulation
définit la manipulation de marché comme incluant les signaux faux ou trompeurs, les négociations à l'ouverture ou à la clôture qui induisent les investisseurs en erreur, les ordres algorithmiques qui
déclenchent ou amplifient une tendance, et la \omterm{def:s2:manipulative-strategies-and-how-they-are-caught:benchmark}{manipulation d'indices de référence}. Son annexe~I énumère des indicateurs, parmi lesquels les transactions sans aucun changement
de propriété effective et les ordres retirés avant exécution après avoir modifié l'apparence du carnet.
\end{dated}

\section{Déclenchement de momentum et manipulation du cours de clôture}

\begin{definition}[Déclenchement de momentum]\label{def:s2:manipulative-strategies-and-how-they-are-caught:ignition}
Le \emph{déclenchement de momentum}\index{déclenchement de momentum} (momentum ignition) consiste à passer des ordres ou des transactions destinés à lancer ou à exagérer une tendance de cours, afin que les réactions des autres intervenants
fassent progresser le cours davantage, et à négocier contre le mouvement qu'ils provoquent.
\end{definition}

\begin{definition}[Manipulation du cours de clôture]\label{def:s2:manipulative-strategies-and-how-they-are-caught:close}
La \emph{manipulation du cours de clôture}\index{manipulation du cours de clôture} (banging the close) consiste à négocier massivement juste avant et pendant la période de clôture afin de déplacer
le cours de clôture ou le prix de règlement-livraison, au profit d'une position dont la valeur dépend de ce prix.
\end{definition}

\textbf{Fonctionnement.} Le Market Abuse Regulation qualifie de manipulation les ordres qui déclenchent ou amplifient une tendance~; le chapitre n'a trouvé aucune affaire de répression
qui utilise l'expression \omterm{def:s2:manipulative-strategies-and-how-they-are-caught:ignition}{déclenchement de momentum}, et la plainte de la CFTC contre Sarao décrit un \omterm{def:s2:manipulative-strategies-and-how-they-are-caught:spoofing}{spoofing} conçu pour provoquer des variations de cours susceptibles d'être exploitées.
La \omterm{def:s2:manipulative-strategies-and-how-they-are-caught:close}{manipulation du cours de clôture} est mieux documentée. Dans l'affaire Optiver menée par la CFTC, des traders ont accumulé de vastes positions Trading at Settlement, valorisées
au prix de règlement du jour, sur les contrats à terme sur pétrole brut, fioul domestique et essence de la NYMEX. Ils ont ensuite négocié des contrats à terme dans le sens opposé juste avant et pendant la
clôture, afin de faire évoluer le règlement dans leur sens.

\textbf{Pourquoi c'est illégal.} La négociation à la clôture qui induit en erreur les investisseurs qui se fondent sur les cours de clôture constitue une manipulation au titre du Market Abuse Regulation,
et la négociation calibrée pour déplacer les prix de règlement et de référence est l'un de ses indicateurs. Aux États-Unis, il s'agit d'une manipulation ou d'une tentative de manipulation
au sens du Commodity Exchange Act.

\textbf{L'affaire de répression.} La plainte fait état de 19 tentatives en mars 2007, dont 5 au moins ont réussi, ainsi que de fausses déclarations à la NYMEX. La surveillance propre
à la NYMEX a détecté les négociations. En avril 2012, une décision d'accord amiable a imposé une sanction civile de 13~millions de dollars et une restitution de 1~million de dollars,
des limites de négociation pour Optiver, ainsi que des interdictions d'exercer de deux à huit ans pour trois anciens dirigeants.

\section{Wash trading et manipulation d'indice de référence}

\begin{definition}[Manipulation d'indice de référence]\label{def:s2:manipulative-strategies-and-how-they-are-caught:benchmark}
La \emph{manipulation d'indice de référence}\index{manipulation d'indice de référence} (benchmark manipulation) consiste à soumettre de fausses données à un indice de référence, ou à négocier ou à se concerter pour modifier les données
à partir desquelles il est calculé, au profit de positions valorisées sur l'indice de référence.
\end{definition}

\textbf{Fonctionnement.} Un wash trade (livre~3, chapitre~16) a le même bénéficiaire effectif des deux côtés, de sorte qu'il crée du volume et des cours sans changement
de propriété. En octobre 2024, la SEC a poursuivi trois sociétés qui se présentaient comme teneurs de marché, ainsi que neuf personnes physiques, pour avoir vendu de la manipulation de marché
comme service à des promoteurs de crypto-actifs. Elles étaient accusées d'avoir négocié entre elles et d'avoir utilisé des robots qui ont parfois généré des volumes artificiels de plusieurs milliards de dollars
par jour. Cong et ses coauteurs ont estimé que le wash trading représentait plus de 70~\% du volume déclaré sur les plateformes de crypto-actifs non régulées. La \omterm{def:s2:manipulative-strategies-and-how-they-are-caught:benchmark}{manipulation d'indice
de référence} est surtout connue à travers les fixings de change. En mai 2015, quatre banques ont accepté de plaider coupable pour complot en vue de manipuler le taux euro-dollar~:
leurs traders, dans un salon de discussion intitulé «~The Cartel~», se sont concertés autour des fixings de la BCE à 13~h~15 et de WM/Reuters à 16~h (livre~2,
chapitre~17).

\textbf{Pourquoi elles sont illégales.} Les transactions sans aucun changement de propriété effective constituent un indicateur de manipulation au titre du Market Abuse Regulation,
et la manipulation du calcul d'un indice de référence relève de son article~12. Aux États-Unis, le wash trading sur contrats à terme est interdit par le
Commodity Exchange Act, et la manipulation de titres par les lois sur les valeurs mobilières.

\textbf{Les affaires de répression.} Les plaidoyers de culpabilité dans l'affaire des changes ont entraîné des amendes de plus de 2.5~milliards de dollars. L'affaire de la SEC en 2024 faisait partie d'une
enquête au cours de laquelle le FBI a créé un jeton qui lui était propre, et les sociétés auraient manipulé le marché de ce jeton.

\section{Attaque sandwich et son statut juridique}

Une attaque sandwich (livre~3, chapitre~22) achète avant le swap en attente d'une victime sur une plateforme décentralisée (DEX) et vend juste après, dans le même bloc,
de sorte que la victime négocie à un moins bon prix. Elle exploite une information que la blockchain rend publique par construction~: les transactions en attente dans un mempool public. Les utilisateurs
peuvent l'éviter en envoyant leurs transactions par des canaux privés, et des services tels que Flashbots Protect dissimulent les transactions aux robots d'attaque sandwich. La
question de savoir si l'attaque sandwich relève du front-running au sens juridique, de la manipulation ou d'aucune des deux reste ouverte~: le chapitre
n'a trouvé aucune décision de justice à ce sujet dans le corpus consulté, et il ne cite aucune affaire.

\section{Surveillance~: détecteurs et erreurs}

\texttt{firm.surveil} construit des jours de compte, une ligne par compte et par jour, à partir de types légitimes choisis pour déjouer les détecteurs naïfs, et injecte 100 épisodes
de chaque manipulation selon les schémas décrits dans les dossiers de répression. Parmi 11\,500 jours de compte légitimes figurent 5\,000 teneurs de marché qui annulent
presque tout, 3\,000 qui renouvellent toutes leurs cotations après chaque exécution, 1\,500 fournisseurs de profondeur du carnet dont les gros ordres éloignés du meilleur cours sont rarement exécutés,
et 2\,000 traders directionnels. Chaque détecteur reçoit un seuil fixé de sorte que 1~\% des jours de compte légitimes soient signalés
(\cref{lst:s2:manipulative-strategies-and-how-they-are-caught:detect})~:

\begin{center}
\small
\begin{tabular}{@{}p{4.2cm}p{4.4cm}cp{4.4cm}@{}}
\toprule
pratique & détecteur & détectés & type légitime le plus signalé\\
\midrule
\omterm{def:s2:manipulative-strategies-and-how-they-are-caught:spoofing}{spoofing} & ratio ordres-transactions & 0~\% & teneurs de marché (1.8~\%)\\
 & écart de taux d'exécution, petits et gros ordres & 37~\% & fournisseurs de profondeur (7.7~\%)\\
 & annulations après exécutions & 51~\% & renouvelleurs de cotations (3.8~\%)\\
 & écart $\times$ annulations & 77~\% & fournisseurs de profondeur (7.7~\%)\\
\omterm{def:s2:manipulative-strategies-and-how-they-are-caught:close}{manipulation du cours de clôture} (marking the close) & part du volume à la clôture & 1~\% & fonds indiciels (3.7~\%)\\
 & part $\times$ retournement du lendemain & 21~\% & fonds indiciels (3.7~\%)\\
 & $\times$ position au règlement-livraison & 61~\% & fonds indiciels (3.3~\%)\\
wash trading & transactions auto-appariées & 74~\% & sociétés multi-algorithmes (4.0~\%)\\
 & part des transactions auto-appariées & 97~\% & sociétés multi-algorithmes (5.0~\%)\\
\bottomrule
\end{tabular}
\end{center}

Le ratio ordres-transactions échoue parce que le comportement qui définit le spoofer n'est pas d'annuler beaucoup~: les teneurs de marché annulent davantage. Il consiste à annuler un type d'ordre
et à en exécuter un autre, ce qui correspond au témoignage de Coscia~: deux tailles d'ordres utilisées, avec des taux d'annulation très différents. Chacune des moitiés de ce schéma
admet des explications innocentes. Les gros ordres des fournisseurs de profondeur sont rarement exécutés~; les renouvelleurs de cotations annulent juste après les exécutions. Ce n'est qu'ensemble que les deux caractéristiques
séparent les spoofers (la \cref{fig:s2:manipulative-strategies-and-how-they-are-caught:roc}). Le volume de la fenêtre de clôture seul signale les fonds indiciels, qui doivent négocier
à la clôture~; l'inversion et une position valorisée au prix de règlement-livraison sont ce qu'ajoute le dossier Optiver. Les auto-appariements signalent des sociétés dont les algorithmes indépendants
se croisent parfois, si bien que la part des transactions d'un compte qui s'auto-apparient fonctionne mieux que leur nombre.

\begin{omfigure}
\begin{tikzpicture}
\begin{axis}[width=11cm,height=5.4cm,xmode=log,xlabel={\small taux de faux positifs (\%, échelle log)},ylabel={\small épisodes de spoofing détectés (\%)},
  tick label style={font=\footnotesize},xmin=0.1,xmax=50,ymin=0,ymax=100,xtick={0.1,1,10},xticklabels={0.1,1,10},grid=major,
  grid style={gray!25},table/col sep=comma,legend style={font=\scriptsize,at={(0.5,-0.3)},anchor=north},legend columns=2]
\addplot[omThm,very thick] table[x=fpr,y=both]{figdata/strategies-2/29-manipulative-strategies-and-how-they-are-caught/roc.csv};
\addplot[omDef,thick,dashed] table[x=fpr,y=cancel]{figdata/strategies-2/29-manipulative-strategies-and-how-they-are-caught/roc.csv};
\addplot[gray,thick] table[x=fpr,y=gap]{figdata/strategies-2/29-manipulative-strategies-and-how-they-are-caught/roc.csv};
\addplot[black,thick,dotted] table[x=fpr,y=otr]{figdata/strategies-2/29-manipulative-strategies-and-how-they-are-caught/roc.csv};
\legend{écart de taux d'exécution $\times$ annulations après exécutions,annulations après exécutions,écart de taux d'exécution,ratio ordres-transactions}
\end{axis}
\end{tikzpicture}

\omcaption{Part des jours de compte de \omterm{def:s2:manipulative-strategies-and-how-they-are-caught:spoofing}{spoofing} injectés détectés par quatre détecteurs, en regard de la part des jours de compte légitimes signalés. Données~:
\texttt{s2\_surveil.roc}.}\label{fig:s2:manipulative-strategies-and-how-they-are-caught:roc}
\end{omfigure}

Une alerte de surveillance marque le point de départ d'une enquête, et non le verdict. Avec un taux de faux positifs de 1~\%, 115 jours de compte légitimes sont signalés pour
chaque 77 épisodes de \omterm{def:s2:manipulative-strategies-and-how-they-are-caught:spoofing}{spoofing} réels détectés dans cet échantillon. La surveillance réelle classe les alertes, lit les ordres qui les sous-tendent et cherche l'intention, comme
le faisaient les tribunaux (la \cref{fig:s2:manipulative-strategies-and-how-they-are-caught:flagged}).

\begin{omfigure}
\begin{tikzpicture}
\begin{axis}[width=11cm,height=5.2cm,ybar,bar width=5pt,ylabel={\small signalés à 1~\% de faux positifs (\%)},tick label style={font=\footnotesize},
  xtick=data,xticklabels={{teneur\\de marché},renouvelleur,{fournisseur\\de profondeur},directionnel,spoofer},x tick label style={align=center},ymin=0,ymax=100,grid=major,grid style={gray!25},
  table/col sep=comma,legend style={font=\scriptsize,at={(0.03,0.97)},anchor=north west}]
\addplot[fill=gray!40,draw=gray] table[x=k,y=otr]{figdata/strategies-2/29-manipulative-strategies-and-how-they-are-caught/flagged.csv};
\addplot[fill=white,draw=black] table[x=k,y=gap]{figdata/strategies-2/29-manipulative-strategies-and-how-they-are-caught/flagged.csv};
\addplot[fill=omDef!50,draw=omDef] table[x=k,y=cancel]{figdata/strategies-2/29-manipulative-strategies-and-how-they-are-caught/flagged.csv};
\addplot[fill=omThm!60,draw=omThm] table[x=k,y=both]{figdata/strategies-2/29-manipulative-strategies-and-how-they-are-caught/flagged.csv};
\legend{ratio ordres-transactions,écart de taux d'exécution,annulations après exécutions,les deux}
\end{axis}
\end{tikzpicture}

\omcaption{Part de chaque type de compte synthétique signalé par les quatre détecteurs de \omterm{def:s2:manipulative-strategies-and-how-they-are-caught:spoofing}{spoofing}, avec des seuils qui signalent 1~\% de l'ensemble des jours de compte légitimes.
Données~: \texttt{s2\_surveil.detectors}.}\label{fig:s2:manipulative-strategies-and-how-they-are-caught:flagged}
\end{omfigure}

\section{Fiches de stratégie}

Les sept pratiques manipulatrices décrites ci-dessous le sont, dans l'ordre qu'exige cette série, selon leur fonctionnement, les raisons pour lesquelles elles sont illégales ou contestées, le
cas de répression et le détecteur de surveillance. La dernière fiche est celle du desk de surveillance lui-même.

\begin{strategyfile}{Spoofing}\label{strat:s2:manipulative-strategies-and-how-they-are-caught:spoofing}
\sfield{Victimes} Les traders qui réagissent aux ordres affichés et négocient à des prix déplacés par des ordres qui n'étaient jamais destinés à être exécutés.
\sfield{Instruments et plateformes} Contrats à terme et autres marchés disposant de carnets d'ordres centraux à cours limité visibles.
\sfield{Fonctionnement} De gros ordres d'un côté, annulés avant de pouvoir être exécutés, et de petits ordres de l'autre côté, exécutés après le déplacement du prix.
\sfield{Pourquoi c'est illégal} Offrir à l'achat ou à la vente dans l'intention d'annuler avant l'exécution~: 7~U.S.C.~6c(a)(5)(C)~; signaux faux ou trompeurs au sens du
Market Abuse Regulation.
\sfield{Cas de répression} United States v. Coscia (7th Cir. 2017)~: 36 mois, condamnation confirmée.
\sfield{Détecteur de surveillance} Comparaison des taux d'exécution des petits et des gros ordres, avec de grandes annulations juste après des exécutions du côté opposé.
\sfield{Détection effective} Surveillance par les bourses et le régulateur des journaux d'ordres~; intention déduite de la conception du programme et des registres de transactions.
\sfield{Évaluer le détecteur en toute honnêteté} Des sosies légitimes dans l'évaluation (fournisseurs de profondeur, renouvelleurs de cotations)~; faux positifs au
seuil utilisé.
\sfield{Sources} Opinion de la Seventh Circuit (2017)~; ce chapitre~: 77~\% détectés à 1~\% de faux positifs.
\end{strategyfile}

\begin{strategyfile}{Layering}\label{strat:s2:manipulative-strategies-and-how-they-are-caught:layering}
\sfield{Victimes} Les traders qui lisent la profondeur du carnet.
\sfield{Instruments et plateformes} Carnets d'ordres profonds et visibles~; le E-mini S\&P 500 dans l'affaire Sarao.
\sfield{Fonctionnement} Plusieurs gros ordres à des niveaux successifs, tenus à l'écart du meilleur prix puis annulés, comme le retrace le CFTC.
\sfield{Pourquoi c'est illégal} \omterm{def:s2:manipulative-strategies-and-how-they-are-caught:spoofing}{Spoofing} au sens du droit américain~; ordres qui modifient la représentation du carnet et sont retirés avant exécution au sens du Market Abuse
Regulation.
\sfield{Cas de répression} CFTC v. Nav Sarao Futures Limited PLC and Sarao (2015)~: sanction et restitution de 38.6~millions de dollars par accord amiable (2016).
\sfield{Détecteur de surveillance} Gros ordres persistants à plusieurs niveaux du meilleur cours, recotés au gré du marché, avec un taux d'exécution quasi nul.
\sfield{Détection effective} Reconstitution du carnet d'ordres~; plaintes d'autres intervenants.
\sfield{Évaluer le détecteur en toute honnêteté} Les teneurs de marché qui cotent de grandes tailles en profondeur se ressemblent tant qu'on n'ajoute pas les transactions du côté opposé.
\sfield{Sources} Communiqués de presse du CFTC 7156-15 et 7486-16.
\end{strategyfile}

\begin{strategyfile}{Déclenchement de momentum (momentum ignition)}\label{strat:s2:manipulative-strategies-and-how-they-are-caught:momentum}
\sfield{Victimes} Les traders dont les algorithmes suivent les tendances de court terme.
\sfield{Instruments et plateformes} Marchés électroniques où interviennent des acteurs de suivi de tendance (trend following).
\sfield{Fonctionnement} Ordres ou transactions destinés à lancer ou à exagérer un mouvement, négociés à contre-sens pendant que d'autres suivent.
\sfield{Pourquoi c'est illégal} Market Abuse Regulation, art.~12(2)(c)(iii)~: ordres qui initient ou accentuent une tendance.
\sfield{Cas de répression} Aucune affaire utilisant ce terme n'a été trouvée~; les affaires de \omterm{def:s2:manipulative-strategies-and-how-they-are-caught:spoofing}{spoofing} décrivent des variations de prix provoquées pour être exploitées.
\sfield{Détecteur de surveillance} Rafales d'ordres agressifs suivies de transactions en sens opposé et d'un retournement (indicateur (e) de l'annexe~I).
\sfield{Détection effective} Alertes sur des schémas de négociation concentrés accompagnés de retournements.
\sfield{Évaluer le détecteur en toute honnêteté} Les algorithmes d'exécution légitimes négocient aussi par rafales~; les retournements sont fréquents.
\sfield{Sources} Market Abuse Regulation, art.~12 et annexe~I.
\end{strategyfile}

\begin{strategyfile}{Wash trades}\label{strat:s2:manipulative-strategies-and-how-they-are-caught:wash}
\sfield{Victimes} Les investisseurs qui lisent le volume et les prix comme un intérêt réel.
\sfield{Instruments et plateformes} Actifs crypto sur des plateformes non réglementées~; contrats à terme~; titres peu liquides.
\sfield{Fonctionnement} Négocier avec soi-même ou avec un compte lié, sans changement de propriétaire effectif.
\sfield{Pourquoi c'est illégal} Interdit sur les marchés américains de contrats à terme et de titres~; indice de manipulation au sens du Market Abuse Regulation.
\sfield{Cas de répression} Accusations de la SEC contre trois prétendus teneurs de marché et neuf personnes physiques (octobre 2024).
\sfield{Détecteur de surveillance} La part des transactions d'un compte appariées avec le même propriétaire effectif.
\sfield{Détection effective} Données de propriété effective entre comptes~; registres de prévention des auto-appariements des bourses.
\sfield{Évaluer le détecteur en toute honnêteté} Des sociétés aux algorithmes indépendants se croisent par accident~; le comptage pénalise les gros négociants honnêtes.
\sfield{Sources} Communiqué de presse de la SEC 2024-166~; Cong et coauteurs (2022)~; ce chapitre~: 97~\% détectés à 1~\% de faux positifs.
\end{strategyfile}

\begin{strategyfile}{Manipulation du cours de clôture}\label{strat:s2:manipulative-strategies-and-how-they-are-caught:close}
\sfield{Victimes} Les détenteurs de positions et de contrats valorisés au cours de clôture ou de règlement.
\sfield{Instruments et plateformes} Règlements de contrats à terme~; enchères de clôture~; contrats Trading at Settlement.
\sfield{Fonctionnement} Une position valorisée au cours de règlement, puis un volume important échangé contre elle dans les dernières minutes de séance, comme dans le compte d'Optiver devant la CFTC.
\sfield{Pourquoi c'est illégal} Manipulation au sens du Commodity Exchange Act~; négociation en clôture qui induit les investisseurs en erreur au sens du Market Abuse
Regulation.
\sfield{Cas de répression} Affaire CFTC c. Optiver (ordonnance par consentement de 2012)~: 14~millions de dollars~; interdictions de négocier pour d'anciens dirigeants.
\sfield{Détecteur de surveillance} Part du volume échangée pendant la fenêtre de clôture, avec un retournement le lendemain et une position valorisée au cours de règlement.
\sfield{Détection effective} Surveillance par les bourses des fenêtres de règlement~; le NYMEX a détecté les transactions d'Optiver.
\sfield{Évaluer le détecteur en toute honnêteté} Les fonds indiciels doivent négocier à la clôture~; le volume seul les signale à tort.
\sfield{Sources} Communiqué de presse de la CFTC 6239-12~; ce chapitre~: 61~\% détectés à 1~\% de faux positifs.
\end{strategyfile}

\begin{strategyfile}{Manipulation d'indice de référence}\label{strat:s2:manipulative-strategies-and-how-they-are-caught:benchmark}
\sfield{Victimes} Toutes les personnes dont les contrats font référence à l'indice.
\sfield{Instruments et plateformes} Fixings de change~; taux d'intérêt de référence~; évaluations de matières premières.
\sfield{Fonctionnement} Fausses soumissions, ou négociation coordonnée autour de la fenêtre de calcul.
\sfield{Pourquoi c'est illégal} Market Abuse Regulation, article 12, paragraphe 1, point d)~; fixation de prix et manipulation au regard du droit américain.
\sfield{Cas de répression} Plaidoyers de culpabilité de quatre banques pour manipulation du fixing de change (mai 2015)~; amendes de plus de 2.5~milliards de dollars.
\sfield{Détecteur de surveillance} Négociation concentrée autour des heures de calcul, avec examen des communications.
\sfield{Détection effective} Surveillance des communications~; contrôles de données par les administrateurs des indices de référence.
\sfield{Évaluer le détecteur en toute honnêteté} Les ordres de fixing des clients sont exécutés légitimement autour du fixing.
\sfield{Sources} Département de la Justice des États-Unis (2015)~; livre~2, chapitre~17.
\end{strategyfile}

\begin{strategyfile}{Attaque sandwich}\label{strat:s2:manipulative-strategies-and-how-they-are-caught:sandwich}
\sfield{Victimes} Les utilisateurs dont les échanges en attente sont visibles dans un mempool public.
\sfield{Instruments et plateformes} Plateformes décentralisées sur les blockchains publiques.
\sfield{Fonctionnement} Un achat avant un échange en attente et une vente après lui, dans le même bloc (livre~3, chapitre~22).
\sfield{Pourquoi c'est contesté} Elle exploite une information rendue publique par conception~; aucune décision de justice à son sujet n'a été trouvée dans les recherches effectuées.
\sfield{Cas de répression} Aucun cas nommé.
\sfield{Détecteur de surveillance} Paires de transactions d'une même adresse encadrant l'échange d'un tiers dans le même pool et le même bloc.
\sfield{Détection effective} On l'évite plutôt qu'on ne la détecte~: les canaux de transaction privés masquent les échanges.
\sfield{Évaluer le détecteur en toute honnêteté} Un arbitrage qui encadre par hasard un échange lui ressemble.
\sfield{Sources} Documentation de Flashbots.
\end{strategyfile}

\begin{strategyfile}{Détecteur~: surveillance des schémas d'annulation}\label{strat:s2:manipulative-strategies-and-how-they-are-caught:detector}
\sfield{Qui paie, et pourquoi} La plateforme, le courtier ou la société elle-même, tenus de détecter et de signaler les ordres suspects.
\sfield{Instruments et plateformes} Tout journal d'ordres que la société ou la plateforme conserve.
\sfield{Signal} Par compte et par jour~: l'écart entre les taux d'exécution des petits et des gros ordres, multiplié par la part des grosses annulations juste après une exécution du
côté opposé.
\sfield{Dimensionnement et exécution} Un seuil fixé au nombre d'alertes que les enquêteurs peuvent lire.
\sfield{Coûts} Temps des enquêteurs~; faux positifs.
\sfield{Comment il échoue} Les manipulateurs qui mélangent tailles et délais~; les stratégies légitimes qui dérivent vers ce schéma.
\sfield{Horizon, capacité, infrastructure} Quotidien~; journaux d'ordres complets et horodatés.
\sfield{Tester le backtest en toute honnêteté} Évaluer sur des imitations légitimes, au seuil opérationnel, avec les épisodes étiquetés à partir de cas réels lorsque c'est possible.
\sfield{Sources} United States v. Coscia (2017)~; ce chapitre~: 77~\% à 1~\% de faux positifs, contre 0~\% pour le ratio ordres-transactions.
\end{strategyfile}

\section{Tutoriel~: dans le journal d'ordres}\label{tut:s2:manipulative-strategies-and-how-they-are-caught:tut}

\begin{tutorial}
\textbf{Objectif.} Construire des comptes-jours de types légitimes et d'épisodes implantés, les noter avec des détecteurs tirés des dossiers de répression, et mesurer
le taux de détection de chaque détecteur à un taux de faux positifs fixé. \textbf{État final~:} le tableau et deux figures.
\begin{tutsteps}
\item \textbf{Les détecteurs de \omterm{def:s2:manipulative-strategies-and-how-they-are-caught:spoofing}{spoofing}}.
\omcode{code/firm/surveil/firm_surveil.py}{67}{76}{Ratio ordres-transactions, écart de taux d'exécution, annulations après exécutions, et les deux.}\label{lst:s2:manipulative-strategies-and-how-they-are-caught:detect}
\item \textbf{Évaluation}.
\omcode{code/firm/surveil/firm_surveil.py}{115}{124}{Le taux de détection au seuil qui signale une part fixe des comptes légitimes.}\label{lst:s2:manipulative-strategies-and-how-they-are-caught:eval}
\item \textbf{Exécuter} \texttt{detectors()}, \texttt{roc()} et \texttt{fig\_surveil.py}.
\end{tutsteps}
\textbf{Pour aller plus loin.} Ajouter un type légitime qui couvre des options au moyen de petits ordres face à de gros cours affichés~; mesurer combien d'alertes par jour un
enquêteur devrait lire~; ajouter les données de position du détecteur de clôture provenant d'une seconde source.
\end{tutorial}

\section{Construction~: surveillance}\label{bld:s2:manipulative-strategies-and-how-they-are-caught:surveil}

\begin{build}
\bfield{Rôle} Comptes-jours étiquetés et détecteurs de \omterm{def:s2:manipulative-strategies-and-how-they-are-caught:spoofing}{spoofing}, de \omterm{def:s2:manipulative-strategies-and-how-they-are-caught:close}{manipulation du cours de clôture} et de wash trades, évalués à des taux de faux positifs fixés.
\bfield{Interface} \texttt{SurveilConfig(\dots)}, \texttt{spoofing\_days}, \texttt{close\_days}, \texttt{wash\_days}, \texttt{spoof\_scores},
\texttt{close\_scores}, \texttt{wash\_scores}, \texttt{tpr\_at\_fpr(score, label, fpr)}, \texttt{flagged\_by\_type}.
\bfield{Règles} Types légitimes choisis pour faire réagir chaque détecteur~; épisodes uniquement tels que décrits dans les dossiers de répression~; aucun profit de manipulation modélisé.
\bfield{Tests d'acceptation} \texttt{code/firm/surveil/tests/}~: le taux de détection à la main~; scores sur deux comptes-jours~; le détecteur combiné dépasse ses parties
et le ratio ordres-transactions échoue~; la part d'auto-appariements à la main.
\bfield{Pour aller plus loin} Journaux au niveau des ordres~; classement des alertes~; détecteurs appris à partir de cas étiquetés.
\end{build}

\begin{omsources}
\item United States v. Coscia, No.~16-3017 (7th Cir. 2017), 866 F.3d 782.
\item Commodity Futures Trading Commission, communiqués de presse 7156-15 (21 avril 2015), 7486-16 (17 novembre 2016) et 6239-12 (19 avril 2012).
\item Règlement (UE) n°~596/2014 sur les abus de marché, article~12 et annexe~I.
\item US Department of Justice, communiqué de presse sur les plaidoyers de culpabilité en change, 20 mai 2015.
\item Securities and Exchange Commission, communiqué de presse 2024-166, 9 octobre 2024.
\item L.~W.~Cong, X.~Li, K.~Tang et Y.~Yang, ``Crypto wash trading'', NBER Working Paper 30783, 2022.
\end{omsources}

\section{Exercices}

\begin{exercise}[$\star$]\label{exo:s2:manipulative-strategies-and-how-they-are-caught:1}
Un trader a passé 24\,814 gros ordres, dont 0.5~\% ont été exécutés. Combien ont été exécutés~?
\end{exercise}

\begin{exercise}[$\star$]\label{exo:s2:manipulative-strategies-and-how-they-are-caught:2}
À un taux de faux positifs de 1~\% sur 11\,500 comptes-jours légitimes, combien sont signalés, et combien d'épisodes réels sur 100 un détecteur au taux de
détection de 77~\% trouve-t-il~?
\end{exercise}

\begin{exercise}[$\star$]\label{exo:s2:manipulative-strategies-and-how-they-are-caught:3}
L'ordonnance par consentement de Sarao exigeait une pénalité de 25.7~millions de dollars et une restitution de 12.9~millions de dollars. Quel est le total~?
\end{exercise}

\begin{exercise}[$\star\star$]\label{exo:s2:manipulative-strategies-and-how-they-are-caught:4}
Pourquoi le ratio ordres-transactions ne parvient-il pas à détecter les spoofers~?
\end{exercise}

\begin{exercise}[$\star\star$]\label{exo:s2:manipulative-strategies-and-how-they-are-caught:5}
Pourquoi le seul volume de la fenêtre de clôture est-il un mauvais détecteur de la \omterm{def:s2:manipulative-strategies-and-how-they-are-caught:close}{manipulation du cours de clôture}~?
\end{exercise}

\begin{exercise}[$\star\star$]\label{exo:s2:manipulative-strategies-and-how-they-are-caught:6}
Pourquoi l'intention, et non l'annulation, doit-elle définir le \omterm{def:s2:manipulative-strategies-and-how-they-are-caught:spoofing}{spoofing}~?
\end{exercise}

\begin{exercise}[$\star\star\star$]\label{exo:s2:manipulative-strategies-and-how-they-are-caught:7}
\emph{Programmation.} Exécuter \texttt{roc()}. Combien d'épisodes de \omterm{def:s2:manipulative-strategies-and-how-they-are-caught:spoofing}{spoofing} le détecteur combiné détecte-t-il pour des taux de faux positifs de 0.1~\% et de 10~\%, et que
cela implique-t-il pour la charge de travail d'un enquêteur~?
\end{exercise}

\begin{exercise}[$\star\star\star$]\label{exo:s2:manipulative-strategies-and-how-they-are-caught:8}
\emph{Trouver la faille.} «~Notre surveillance n'a signalé personne dont le ratio ordres-transactions dépassait notre limite cette année, donc il n'y a pas de \omterm{def:s2:manipulative-strategies-and-how-they-are-caught:spoofing}{spoofing} chez nous.~»
\end{exercise}

\section{Problème~: dans le journal des ordres}

\begin{problem}[{Problème du week-end --- stratégies manipulatrices et surveillance}]\label{pb:s2:manipulative-strategies-and-how-they-are-caught:1}
Le registre public des sanctions et les journées de compte synthétiques du chapitre.

\medskip\noindent\textbf{Partie I --- \omterm{def:s2:manipulative-strategies-and-how-they-are-caught:spoofing}{Spoofing}.}
\begin{enumerate}
\item Définir le \omterm{def:s2:manipulative-strategies-and-how-they-are-caught:spoofing}{spoofing} et le \omterm{def:s2:manipulative-strategies-and-how-they-are-caught:layering}{layering}.
\item Qu'a montré le relevé d'ordres de Coscia~?
\item Qu'a allégué la CFTC au sujet de Sarao~?
\item Qu'interdit le droit américain, et que recense la Market Abuse Regulation~?
\end{enumerate}

\medskip\noindent\textbf{Partie II --- Autres pratiques.}
\begin{enumerate}[resume]
\item Définir le \omterm{def:s2:manipulative-strategies-and-how-they-are-caught:ignition}{déclenchement de momentum} (momentum ignition) et la \omterm{def:s2:manipulative-strategies-and-how-they-are-caught:close}{manipulation du cours de clôture} (marking the close).
\item Que s'est-il passé dans l'affaire Optiver~?
\item Définir la \omterm{def:s2:manipulative-strategies-and-how-they-are-caught:benchmark}{manipulation d'indice de référence}~; que s'est-il passé lors des fixings de change~?
\item Qu'a allégué la SEC en 2024~?
\end{enumerate}

\medskip\noindent\textbf{Partie III --- Surveillance.}
\begin{enumerate}[resume]
\item Décrire les types de comptes synthétiques et la raison d'être de chacun.
\item Donner le tableau des détecteurs.
\item Pourquoi la combinaison de deux caractéristiques fonctionne-t-elle~?
\item Quel est le statut juridique du sandwiching~?
\end{enumerate}

\medskip\noindent\textbf{Partie IV --- Le verdict.}
\begin{enumerate}[resume]
\item Énoncer le \emph{résultat nommé}~: le taux de vrais positifs de chaque détecteur à un taux de faux positifs fixé, et la manière dont la teneur de marché légitime les déclenche.
\item Que se passe-t-il pour une alerte après le déclenchement du détecteur~?
\item Comment évaluer honnêtement un détecteur~?
\item Que modifierait un manipulateur pour échapper au détecteur combiné, et quel en serait le coût pour la manipulation~?
\item Pourquoi les affaires de sanction importent-elles pour la conception des détecteurs~?
\item Quelle fiche de stratégie est légitime~?
\item Quel est le lien entre ce chapitre et les mark-outs du chapitre~24~?
\item En une phrase~: qu'est-ce qui distingue une teneur de marché d'un spoofer~?
\end{enumerate}
\end{problem}

\section{Questions d'entretien}

\begin{interviewq}[$\star$ \iqroles{trader}]\label{iq:s2:manipulative-strategies-and-how-they-are-caught:1}
Vous annulez 95~\% de vos ordres. Comment montrez-vous que vous ne faites pas de \omterm{def:s2:manipulative-strategies-and-how-they-are-caught:spoofing}{spoofing}~?
\end{interviewq}

\begin{interviewq}[$\star\star$ \iqroles{chercheur}]\label{iq:s2:manipulative-strategies-and-how-they-are-caught:2}
Concevez un détecteur de \omterm{def:s2:manipulative-strategies-and-how-they-are-caught:spoofing}{spoofing} à partir d'un journal d'ordres, et dites comment vous le testeriez.
\end{interviewq}

\begin{interviewq}[$\star\star$ \iqroles{trader}]\label{iq:s2:manipulative-strategies-and-how-they-are-caught:3}
Vous détenez une position importante valorisée au règlement-livraison du jour. Que pouvez-vous faire dans les dernières minutes de la séance~?
\end{interviewq}

\begin{interviewq}[$\star\star$ \iqroles{risque}]\label{iq:s2:manipulative-strategies-and-how-they-are-caught:4}
Quels contrôles une société de trading doit-elle mettre en place contre ses propres algorithmes qui manipuleraient les marchés~?
\end{interviewq}

\begin{interviewq}[$\star\star$ \iqroles{développeur}]\label{iq:s2:manipulative-strategies-and-how-they-are-caught:5}
Que doit enregistrer un journal d'ordres pour que la surveillance fonctionne~?
\end{interviewq}

\begin{interviewq}[$\star\star\star$ \iqroles{chercheur}]\label{iq:s2:manipulative-strategies-and-how-they-are-caught:6}
Un détecteur signale 1~\% des 11\,500 journées de compte légitimes et 77 épisodes sur 100. Quelle part des alertes sont vraies, et que devient-elle si les
épisodes sont dix fois plus rares~?
\end{interviewq}
